\documentclass[runningheads]{llncs}

% TODO REVIEW: Insert your submission number below by replacing '*****'
% TODO FINAL: Comment out the following line for the camera-ready version
\usepackage[review,year=2026,ID=*****]{eccv}
% TODO FINAL: Un-comment the following line for the camera-ready version
%\usepackage{eccv}

\usepackage{eccvabbrv}
\usepackage{graphicx}
\usepackage{booktabs}
\usepackage{placeins}

\usepackage[accsupp]{axessibility}

% TODO FINAL: Comment out the following line for the camera-ready version
%\usepackage[pagebackref,breaklinks,colorlinks,citecolor=eccvblue]{hyperref}
% TODO FINAL: Un-comment the following line for the camera-ready version
\usepackage{hyperref}

% Supplementary numbering convention
\renewcommand{\thesection}{\Alph{section}}
\renewcommand{\thetable}{S\arabic{table}}
\renewcommand{\thefigure}{S\arabic{figure}}

\begin{document}

\title{Supplementary Material:\\RefCompose: Multi-Reference Image Generation via LoRA-Conditioned Diffusion}

\titlerunning{Supplementary Material: RefCompose}

\author{First Author\inst{1} \and
Second Author\inst{2,3} \and
Third Author\inst{3}}

\authorrunning{F.~Author et al.}

\institute{Princeton University, Princeton NJ 08544, USA \and
Springer Heidelberg, Tiergartenstr.~17, 69121 Heidelberg, Germany \and
ABC Institute, Rupert-Karls-University Heidelberg, Heidelberg, Germany}

\maketitle

% ------------------------------------------------------------------
\section{Overview}
\label{sec:overview}
% ------------------------------------------------------------------

This supplementary document accompanies the main ECCV workshop paper on \textbf{RefCompose}.
The main paper focuses on Dense Layout quantitative evaluation on InstanceAssemble~\cite{xiang2025instanceassemble}; here we provide:
\begin{enumerate}
  \item \textbf{Layout-centric benchmarks} (Section~\ref{sec:layout_eval}) comparing the prompt-induced spatial layout stage against dedicated layout generators on COCO-GR~\cite{feng2024layoutgpt} numerical and spatial reasoning splits.
  \item \textbf{Extended ablation study} (Section~\ref{sec:ablation}) with quantitative metrics complementing the qualitative analysis in the main paper.
  \item \textbf{Training data details} (Section~\ref{sec:dataset}) referenced from the main experiments section.
\end{enumerate}
Throughout, multiple visual references are composed under a single scene prompt without annotated bounding boxes or per-reference token scaling.

% ------------------------------------------------------------------
\section{Layout-Centric Evaluation}
\label{sec:layout_eval}
% ------------------------------------------------------------------

We situate the spatial layout stage of RefCompose in the broader layout-generation literature by comparing against LayoutGPT~\cite{feng2024layoutgpt} and Lay-Your-Scene~\cite{srivastava2025layyourscene} on COCO-GR~\cite{feng2024layoutgpt} numerical and spatial reasoning splits.
RefCompose first induces a coarse composition with a frozen text-to-image model (Flux~\cite{esser2024flux}) and then extracts object boxes with Grounding DINO~\cite{liu2023grounding}.

\paragraph{Metrics.}
\textbf{L-FID} evaluates layout realism by rendering predicted layouts as colour-coded object maps and computing an FID-style distance to ground-truth layouts~\cite{heusel2017fid}; lower is better.
For \emph{numerical reasoning}, precision, recall, and accuracy measure whether the predicted object set matches the prompt-specified objects.
\textbf{Numerical GLIP-Accuracy} checks whether the requested object counts appear in the generated image~\cite{li2022glip}.
For \emph{spatial reasoning}, accuracy measures whether predicted boxes satisfy relations such as \emph{left}, \emph{right}, \emph{above}, or \emph{below}; \textbf{spatial GLIP-Accuracy} checks whether those relations remain visible after image generation~\cite{li2022glip}.

\paragraph{Results.}
Table~\ref{tab:layout_lfid} reports layout quality; Table~\ref{tab:layout_reasoning} reports numerical and spatial reasoning.
RefCompose attains the best L-FID and spatial accuracy among compared methods.
We attribute this to inverse layout construction: a large pretrained text-to-image model~\cite{esser2024flux} first produces an aesthetically composed scene, after which boxes are recovered from visible object support.
Unlike LLM planners that reason symbolically~\cite{feng2024layoutgpt}, or layout diffusion transformers trained on narrower layout corpora~\cite{srivastava2025layyourscene}, the T2I prior encodes broad composition statistics from large-scale pretraining; sharper extracted boxes then improve downstream canvas placement and image-level spatial compliance.

\begin{table}[t]
  \centering
  \caption{Layout quality on COCO-GR~\cite{feng2024layoutgpt} (lower L-FID is better).}
  \label{tab:layout_lfid}
  \setlength{\tabcolsep}{8pt}
  \begin{tabular}{lc}
    \toprule
    \textbf{Method} & \textbf{L-FID}$\downarrow$ \\
    \midrule
    LayoutGPT~\cite{feng2024layoutgpt} & 3.51 \\
    Lay-Your-Scene~\cite{srivastava2025layyourscene} & 3.07 \\
    RefCompose (spatial layout) & \textbf{2.98} \\
    \bottomrule
  \end{tabular}
\end{table}

\begin{table}[t]
  \centering
  \caption{Numerical and spatial reasoning on COCO-GR~\cite{feng2024layoutgpt}.
  GLIP-Accuracy measures image-level count/spatial compliance~\cite{li2022glip}.}
  \label{tab:layout_reasoning}
  \resizebox{\linewidth}{!}{%
  \setlength{\tabcolsep}{6pt}
  \begin{tabular}{lcccccc}
    \toprule
    & \multicolumn{4}{c}{\textbf{Numerical Reasoning}}
    & \multicolumn{2}{c}{\textbf{Spatial Reasoning}} \\
    \cmidrule(lr){2-5} \cmidrule(lr){6-7}
    \textbf{Method}
    & Prec.$\uparrow$ & Recall$\uparrow$ & Acc.$\uparrow$ & GLIP$\uparrow$
    & Acc.$\uparrow$ & GLIP$\uparrow$ \\
    \midrule
    LayoutGPT~\cite{feng2024layoutgpt}
      & 76.29 & 86.64 & 76.72 & 54.25 & 87.07 & 56.89 \\
    Lay-Your-Scene~\cite{srivastava2025layyourscene}
      & 77.62 & 99.23 & 95.14 & 56.20 & 92.58 & 58.94 \\
    RefCompose (spatial layout)
      & \textbf{78.21} & \textbf{99.25} & \textbf{96.41} & \textbf{56.29} & \textbf{93.13} & \textbf{59.39} \\
    \bottomrule
  \end{tabular}%
  }
\end{table}

% ------------------------------------------------------------------
\section{Extended Ablation Study}
\label{sec:ablation}
% ------------------------------------------------------------------

The main paper presents a qualitative ablation of the reference canvas and depth prior (Figure~7 in the main document).
Table~\ref{tab:ablation} reports the corresponding quantitative results on the Dense Layout test split~\cite{xiang2025instanceassemble}, using the same evaluation protocol as Table~2 in the main paper.

\begin{table}[t]
  \centering
  \caption{Ablation study evaluating the contributions of the reference canvas and depth prior.
  The full model consistently achieves the best performance across identity preservation, spatial arrangement, and semantic alignment metrics.}
  \label{tab:ablation}
  \resizebox{\linewidth}{!}{%
  \setlength{\tabcolsep}{5pt}
  \begin{tabular}{lccccccccc}
    \toprule
    \textbf{Configuration}
    & \textbf{Spatial}$\uparrow$
    & \textbf{Color}$\uparrow$
    & \textbf{Texture}$\uparrow$
    & \textbf{Shape}$\uparrow$
    & \textbf{PICK}$\uparrow$
    & \textbf{DINO}$\uparrow$
    & \textbf{DINO$_{\text{refs}}$}$\uparrow$
    & \textbf{CLIP-I}$\uparrow$
    & \textbf{CLIP-T}$\uparrow$ \\
    \midrule
    Full model (canvas + depth, dual LoRA)
      & \textbf{0.908} & \textbf{0.826} & \textbf{0.854} & \textbf{0.849}
      & \textbf{0.229} & \textbf{0.959} & \textbf{0.613} & \textbf{0.766} & 0.327 \\
    w/o reference canvas (text-only appearance)
      & 0.861 & 0.542 & 0.571 & 0.558
      & 0.224 & 0.712 & 0.471 & 0.702 & \textbf{0.334} \\
    w/o depth prior (canvas only, no structure)
      & 0.683 & 0.781 & 0.809 & 0.612
      & 0.221 & 0.798 & 0.587 & 0.741 & 0.310 \\
    \bottomrule
  \end{tabular}%
  }
\end{table}

\paragraph{Discussion.}
Removing the \emph{reference canvas} causes the largest drop on Color, Texture, DINO$_{\text{refs}}$, and CLIP-I---appearance-centric signals---while Spatial remains relatively high because the depth prior still constrains placement.
The model then infers appearance from text alone, landing near the CreatiLayout/InstanceAssemble range on photometric metrics.
CLIP-T increases slightly, mirroring the text-alignment trade-off discussed in the main paper.

Removing the \emph{depth prior} causes the largest drop on Spatial and Shape, reflecting degraded placement and occlusion reasoning, while Color, Texture, and DINO$_{\text{refs}}$ remain closer to the full model since the canvas still supplies appearance.
This matches the qualitative observation in the main paper that subjects such as the bat and bed remain recognizable, but their arrangement becomes contextually implausible.
Together, these results confirm that canvas and depth provide complementary, disentangled signals for appearance and structure.

% ------------------------------------------------------------------
\section{Training Data Details}
\label{sec:dataset}
% ------------------------------------------------------------------

As noted in the main paper, RefCompose is trained on a mixture of proprietary cinematic imagery and synthetic composites.

\paragraph{Real data.}
We curate ultra high-resolution movie frames at $1280 \times 720$ from a proprietary catalogue.
Object references are canonicalized to neutral pose and lighting to reduce spurious correlations between scene context and subject appearance.
The five-stage pipeline described in the main paper (Qwen3-VL~\cite{qwen3vl} enumeration, Grounding DINO~\cite{liu2023grounding} localisation, Flux~\cite{esser2024flux} canonicalization, DINO filtering, and re-annotation) yields approximately $50{,}000$ samples from $1{,}000$ distinct films.

\paragraph{Synthetic data.}
We complement real frames with synthetic composites generated using Flux~2.0~\cite{esser2024flux}.
Synthetic scenes emphasise diverse object positions, scales, and orientations that are underrepresented in canonicalized movie crops, while real data supplies richer textures and identity cues.
Each synthetic sample follows the same annotation protocol as real data: detected crops, neutral-pose canonicalization, DINO similarity filtering ($>0.90$), and dual-level Qwen3-VL~\cite{qwen3vl} captions.
This mixture encourages the model to generalise spatial composition beyond studio lighting and fixed poses without sacrificing photometric fidelity on in-domain cinematic content.

\paragraph{Training schedule.}
LoRA training proceeds in three stages over $50{,}000$ iterations at $1280 \times 720$, as summarised in the main paper:
(1)~depth-only warm-up on a black canvas ($10{,}000$ iterations);
(2)~appearance grounding with crops placed at Grounding DINO boxes on ground-truth images ($20{,}000$ iterations);
(3)~robustness fine-tuning with photometric jitter and mild spatial perturbations (remaining iterations).
Rank $r{=}32$ low-rank adapters are injected into double-stream attention blocks following EasyControl~\cite{zhang2025easycontrol}.

\FloatBarrier

\bibliographystyle{splncs04}
\bibliography{main}

\end{document}
